\section*{Task 2: Inversion formula}

\paragraph{Goal.}
Recover the multiplicity coordinates from the lifted rank pattern.

\paragraph{Boundary convention.}
If \((u,v)\notin \widetilde{\mathcal I}_L^{\mathrm{wt}}\), define
\[
\rho_{uv}:=0.
\]
With this convention, every finite-difference formula may be written without
separate boundary cases.

\paragraph{Inverse formula.}
For every admissible lifted interval \((a,b)\in \widetilde{\mathcal I}_L^{\mathrm{wt}}\),
define
\[
\mu_{ab}(\rho)
:=
\rho_{ab}-\rho_{a-1,b}-\rho_{a,b+1}+\rho_{a-1,b+1}.
\]
Then for the rank pattern coming from a multiplicity pattern \(m\), one has
\[
m_{ab}=\mu_{ab}(\rho).
\]

\paragraph{Proof sketch.}
Insert the forward formula into the right-hand side:
\[
\mu_{ab}(\rho)
=
\sum_{u\le a\le b\le v} m_{uv}
-
\sum_{u\le a-1\le b\le v} m_{uv}
-
\sum_{u\le a\le b+1\le v} m_{uv}
+
\sum_{u\le a-1\le b+1\le v} m_{uv}.
\]
For a fixed interval \((u,v)\), the coefficient of \(m_{uv}\) is
\[
\mathbf 1_{u\le a\le b\le v}
-
\mathbf 1_{u\le a-1\le b\le v}
-
\mathbf 1_{u\le a\le b+1\le v}
+
\mathbf 1_{u\le a-1\le b+1\le v}.
\]
This coefficient is \(1\) exactly when
\[
u=a,\qquad v=b,
\]
and it is \(0\) otherwise. Therefore only \(m_{ab}\) survives.

\paragraph{Special cases.}
For the unique maximal interval \((0,L)\), the boundary convention gives
\[
m_{0L}=\rho_{0L},
\]
because all three correction terms vanish. This is the clean lifted expression
of the fact that the end-to-end rank equals the multiplicity of the unique
length-\(L\) indecomposable.

\paragraph{Status.}
Completed. The lifted finite-difference formula recovers all multiplicities and
fixes the main combinatorial defect of the cyclic indexing used in the older
draft.
